% ============================================================
% NEW §III-D: Relation to Prior Time-Series Evaluation Frameworks
% Insert as new subsection after §III-C (Window-Based Instantiation)
% ============================================================

\subsection{Relation to Prior Time-Series Evaluation Frameworks}
\label{sec:related_eval}

Several evaluation frameworks for time-series anomaly detection have been proposed;
we position our work by identifying the axes on which our contributions are distinct.

\begin{table}[!t]
\caption{Comparison of time-series evaluation frameworks.
\emph{Causality}: metric is computed without future information;
\emph{Event agg.}: consecutive abnormal windows are aggregated into events;
\emph{DR@$\Delta t$}: explicit time-budgeted detection rate;
\emph{Deployment semantics}: thresholds are grounded in physical operating constraints.}
\label{tab:eval_comparison}
\centering
\small
\begin{tabular}{lcccc}
\hline
\textbf{Framework} & \textbf{Causality} & \textbf{Event agg.} & \textbf{DR@$\Delta t$} & \textbf{Deploy. sem.} \\
\hline
Tatbul et al.\ \cite{tatbul2018precision} & post-hoc & segment-level & \texttimes & \texttimes \\
NAB Score~\cite{lavin2015evaluating}      & reward-shaping & window-weighted & \texttimes & \texttimes \\
Classical IDS metrics~\cite{axelsson2000base} & online & \texttimes & \texttimes & partial \\
\textbf{This work} & \textbf{online} & \textbf{FP events} & \textbf{\checkmark} & \textbf{risk-based} \\
\hline
\end{tabular}
\end{table}

\textbf{Tatbul et al.'s point-adjusted (PA) and range-based metrics}
\cite{tatbul2018precision} score anomaly \emph{segments} against predicted ranges.
Their metrics are computed in a post-hoc, offline manner and reward any detection
within a labelled anomaly range equally regardless of when the first detection occurs.
Our DR@$\Delta t$ instead scores only the \emph{first detection event} within a
specified time budget, enforcing causality and response-time accountability.

\textbf{The Numenta Anomaly Benchmark (NAB) score} \cite{lavin2015evaluating}
applies a window-weighted reward that discounts late detections and penalises
false positives via score penalties. While NAB is causal, it does not expose
the time-to-first-detection as an interpretable parameter, making it difficult
to translate a score into a concrete operational response-time budget.

\textbf{Classical IDS false-alarm metrics} (FAR, FPR) quantify false positives
as a fraction of normal-class windows \cite{axelsson2000base}. These metrics
treat each window as independent and therefore count the same persistent alarm
multiple times. Our MTBFA instead aggregates consecutive false-positive windows
into a single event—matching operator experience of a false alarm as a single
disruptive episode—and expresses the rate in operational time (hours), a unit
directly interpretable by flight operators.

In summary, compared with existing frameworks our evaluation makes
three specific additions: (i) an explicit, user-specified time budget for
detection recall (DR@$\Delta t$); (ii) event-level (rather than window-level)
false-alarm counting for MTBFA; and (iii) threshold definitions grounded in
UAV-specific physical risk analysis (Section~\ref{sec:deployability_standard}).
